\section{چرا این مسئله مهم است}

\begin{frame}{از هر ۱۰۰ پرس، ۸ تا دور ریخته می‌شود}
  \begin{columns}
    \begin{column}{0.42\textwidth}
      \begin{itemize}
        \item $8.05\%$ رزروها تحویل گرفته نمی‌شوند
        \item سالانه حدود ۲۳۲ هزار پرس، حدود ۸۷ میلیارد تومان به قیمت شهریور ۱۴۰۵
      \end{itemize}
    \end{column}
    \begin{column}{0.55\textwidth}
      \begin{figure}
        \centering
        \includegraphics[width=0.92\textwidth]{slides_01_waffle}
      \end{figure}
    \end{column}
  \end{columns}
  \note{این عدد وزنی است نه میانگین ساده. میانگین ساده $9.59\%$ می‌شود و به سلف‌های کوچکِ
  پرنوسان وزن بیش‌ازحد می‌دهد. اختلاف $19\%$ نسبی است. شاهد: \lr{F01} در data\_facts\_register.}
\end{frame}

\begin{frame}{گره تصمیم‌گیری: نامتقارنی هزینه‌ها}
  \begin{itemize}
    \setlength{\itemsep}{18pt}
    \item \textbf{پخت روی سقف رزرو:} ریسک اعتراض به صفر می‌رسد، اما هدررفت بودجه تضمین‌شده است
    \item \textbf{هزینه‌ی دو خطا برابر نیست:} صدای کمبود غذا بلند و فوری است، ولی اتلاف میلیاردی پنهان و بی‌صدا
    \item \textbf{کم‌پختن بدون ضابطه در عمل:} پیمانکار تجربی کمتر می‌پزد ولی پول کامل می‌گیرد؛ مدل این تصمیم را شفاف و بهینه می‌کند
  \end{itemize}
  \note{نکات کلیدی برای ارائه:
  ۱. پارادوکس مدیریتی: چرا مدیر سلف نمی‌تواند به‌سادگی بگوید «۵٪ کمتر بپزیم»؟ چون اگر یک روز غذا کم بیاید، اعتراض صنفی فوراً شخص مدیر را هدف می‌گیرد؛ اما اگر غذا اضافه بیاید، هزینه در بودجه‌ی عمومی مستهلک می‌شود و هزینه‌ی سیاسی ندارد. محافظه‌کاری در این ساختار یک رفتار کاملاً عقلانی است.
  ۲. واقعیت میدانی (شاهد: project-definition.md بند ۱-۱): مصاحبه‌ها نشان داد بعضی پیمانکاران تجربی همین حالا هم کمتر از رزرو می‌پزند اما فاکتور کامل می‌گیرند؛ یعنی کم‌پختن بدون ضابطه در حال رخ دادن است، اما نفع اقتصادی آن به دانشگاه نمی‌رسد و ریسک کنترل‌نشده ایجاد می‌کند.
  ۳. مأموریت مدل: تبدیل این خلأ مدیریتی به یک تصمیم بهینه، داده‌محور و قابل‌حسابرسی با موازنه‌ی ریاضی میان هزینه‌ی کمبود ($C_u$) و هزینه‌ی مازاد ($C_o$).
  ۴. این اسلاید عمداً تنها اسلاید بدون نمودار ارائه در این بخش است تا گره‌ی اقتصادی و انگیزشی مسئله پیش از ورود به مبانی ریاضی به‌روشنی در ذهن مخاطب تثبیت شود.}
\end{frame}
